\section{沿研究脉络设计跨模型比较，而不是只复制一条率曲线}
第3章重排后的研究线索要求数值设置也相互衔接。下表列出本批全文可确认的输入和不能擅自补全的缺口；它们是文献参数，不是本报告声称已经重跑的算例。
\begin{longtable}{P{29mm}P{119mm}}
\toprule 对照环节 & 可复查的计算条件及限制\\\midrule\endfirsthead
\toprule 对照环节 & 可复查的计算条件及限制\\\midrule\endhead
有限长X线：F659 &
PIC域$32\times64\times16d_i$，$768\times1536\times384$格，质量比75；初始薄/厚半宽$0.5/4d_i$，厚度过渡尺度$5d_i$。薄片长度取31、12、8.4、$4d_i$。$x,y$周期，$z$导体及粒子反射。主索引为9页预印本；不能从这组参数宣称真实质量比收敛。\\
有限长X线：F792 &
同类几何，VPIC为$640\times1280\times320$格、质量比25、$c=20\VA$、上游$\beta=0.3$，每物种超过$2.6\times10^{10}$粒子。八个薄片长度为5--$40d_i$。除固定长度外，还须输出局地$d_e$、片半厚、$L_a(t)$及两种归一化率。\\
各向异性闭合：F611 &
混合与PIC都用$40\times20d_i$域；前者$1024^2$格，后者$5120\times2560$格、质量比400、每物种约400粒/格。比较force-free和Harris初值，导场比0.3及0.2。正文给出$\eta,\eta_H$的方程位置却未列具体系数，重现前必须补齐，不能凭相邻论文猜测。\\
嵌套四极：L213 &
EMHD以$d_e,\Omega_{ce}^{-1}$归一，关键图时刻45和145；主/次X点强度不同。本文未完整给网格、时间步与收敛信息，2009年前序论文目前只补登记待读。虚拟轨道和坐标变换也是模型—观测比较的输入。\\
加热候选：L272 &
Boris试粒子分别1000和10000条轨道；PIC质量比100，初始半厚$d_i$，导场比0、0.1、0.3、0.5。PIC图复用2019前序结果，域、网格、粒子数仍需追溯原方法。不得把未取得的输入写成已复现配置。\\
\bottomrule
\end{longtable}
\citep{F659,F792,F611,L213,L272}

\paragraph{建议比较顺序。}
先在同一初值中确认各代码解对各自方程，再检查同一通量阶段的局地场、电流层和压力预算。对于有限长X线，第一组比较应保持薄/厚片过渡、种子及边界一致，输出整个X线和活跃区两种率；随后才扫描电子尺度与闭合。对于F611类闭合，至少同时比较率、长层长度、$F_e$、面外磁场剖面及Harris参考密度修正。把不同种子的曲线平移到重合只能检验相应演化阶段，不能证明启动时钟已经一致。

\paragraph{加热诊断须单独收敛。}
给定场轨道先检查步长、能量及轨道分离对初始间距的依赖；自洽计算再输出中心速度二阶矩、$\vct P_s:\grad\vct u_s$、热流和粒子能量增量。场梯度指标涉及空间导数，往往比磁场本身更容易受网格噪声影响。比较实验还须对模拟作同等空间、时间和视线平均，不能以更尖的裸网格峰值宣称比实验更强。
